\documentclass[11pt]{article}
\usepackage[T1]{fontenc}
\usepackage{lmodern,amsmath,amssymb,amsthm,mathtools,microtype}
\usepackage[margin=1in]{geometry}
\usepackage[hidelinks]{hyperref}
\newtheorem{theorem}{Proposed theorem}
\newtheorem{lemma}{Lemma}
\newtheorem{proposition}{Proposition}
\newcommand{\ReH}{\operatorname{Re}_{H}}
\newcommand{\ImH}{\operatorname{Im}_{H}}
\newcommand{\norm}[1]{\left\lVert#1\right\rVert}
\newcommand{\abs}[1]{\left|#1\right|}
\newcommand{\ip}[2]{\left\langle#1,#2\right\rangle}
\newcommand{\dd}{\,\mathrm d}
\newcommand{\ee}{\mathrm e}
\title{The exact-floor boundary in broadband absorption asymmetry\\
\large A passive narrow-resonance construction}
\author{Research working note}
\date{17 October 2026}
\begin{document}
\maketitle
\begin{abstract}
We propose a constructive resolution of the boundary condition
$E_{\min}=E_*$ in the resource-unrestricted two-layer tensor class.
A single compact body has exactly prescribed extinction at a chosen interior
frequency, that frequency is its strict global extinction minimum on the
prescribed band, and its absorption-asymmetry ratio is uniformly arbitrarily
close to one. The new ingredient is a small passive Lorentz resonance with a
negative extinction-interference lobe. We first align a minimum of a fixed
planar reference spectrum, then choose a finite aperture, and only afterwards
narrow the added resonance. Material-parameter analyticity transfers the
shape of this lobe to the finite body; ordinary smoothness of that selected
body suffices to place its minimum exactly. No frequency-derivative aperture
limit is assumed. This is an analytical research draft building on the earlier
fixed-band note, not a new formal verification or external peer review.
\end{abstract}

\section{Statement and dependencies}
Let $B=[\omega_1,\omega_2]$, with
$0<\omega_1<\omega_c<\omega_2$, and let $k=\omega/c_{\rm v}$.
Use the two opposite real linearly polarized incident waves
$f_\pm=e_y\ee^{\pm ikz}$.
The bodies have two homogeneous layers, a common finite square aperture,
symmetric local dielectric tensors, vacuum exterior and $\mu=I$.
As in \cite{base,strict}, reciprocal extinction and power balance give
\[
 X_+=X_-=X=A_\pm+S_\pm,\qquad
 \Gamma=\frac{\abs{A_+-A_-}}{X}.
\]

\begin{theorem}[Exact extinction floor at an interior center]\label{thm:main}
Fix $P>1$, $Q>0$, $E_*>0$ and $0<\epsilon<1$.
There is a finite two-layer body with one fixed passive causal reciprocal
dielectric law in each layer such that, on $B$,
\begin{gather}
 I\preceq\ReH\varepsilon\preceq PI,\qquad
 0\prec\ImH\varepsilon\preceq QI,\label{eq:caps}\\
 X(\omega_c)=E_*,\qquad
 X(\omega)>E_*\quad(\omega\ne\omega_c),\qquad
 1-\epsilon<\Gamma(\omega)<1.\label{eq:target}
\end{gather}
The selected susceptibilities $\chi=\varepsilon-I$ are finite sums of damped
Lorentz oscillators with real
positive semidefinite residues and strictly positive damping.
In particular, for the class satisfying $X(\omega_c)=E_*=\min_B X$,
\[
 \sup_D\inf_{\omega\in B}\Gamma_D(\omega)=1,
\]
and the supremum is not attained.
\end{theorem}

The same qualifications as in \cite{strict} apply: caps are imposed on the
specified band, and there is no common bound on aperture, thickness, volume,
oscillator strength or count, or a positive lower bound on material loss or
resonance linewidth. All selected parameters are finite and every damping is
positive. The final spectrum has a dip at the center; it is not constant.

We inherit the weak-Maxwell volume representation, radiation identity,
optical theorem and bounded-domain operator framework of \cite{base}.
We also use the normalized matrix factorization and the $C^0$ frequency-uniform
aperture theorem of \cite{strict}. The additions proved below are an extension
to a complex material parameter, a nonzero material sensitivity, and a
finite-body narrow-resonance minimum argument. Neither the earlier Lean
project nor the numerical checks formally verify this extension.

\section{Why the straightforward repairs are insufficient}
Uniform approximation to a flat or convex spectrum does not control its
derivatives. Small ripples can create competing minima. Maxwell scaling of a
\emph{finite} body obeys
\[
 D\mapsto\lambda D,\qquad
 \chi_{\rm new}(\omega)=\chi_{\rm old}(\lambda\omega),\qquad
 X_{\rm new}(\omega)=\lambda^2X_{\rm old}(\lambda\omega).
\]
Selecting a minimizer $\omega_m$ and setting $\lambda=\omega_m/\omega_c$
places a minimum at the center but changes its value. If the selected
minimizer jumps during amplitude tuning, the resulting
$(\omega_m/\omega_c)^2\min X$ can jump too; a scalar intermediate-value
argument then fails.

Nor may one simply differentiate the old aperture proof twice.
The absolute Weyl-kernel estimate has an integrable
$\abs{\rho-k}^{-1/2}$ grazing singularity; differentiating introduces
nonintegrable absolute majorants. Even the weak thick planar limit need not
converge in $C^2$: an $O(N^{-2})$ Fabry--Perot term oscillating at frequency
$O(N)$ has a second derivative of order one. These observations invalidate
those shortcuts, not the boundary statement itself.

Our scaling is performed on a single \emph{planar density}, once, before
any finite aperture is selected. The later sharp minimum is controlled by a
material parameter rather than an aperture-uniform frequency derivative.

\section{Freeze a planar minimum and a nonzero sensitivity}
We first work on a slightly enlarged compact band $J\Subset(0,\infty)$.
Choose an interval $U$ about $\omega_c$ so small that
\[
 \frac{\nu}{\omega_c}B\subset J\quad\text{for every }\nu\in U.
\]
In the shaping construction of \cite{strict}, replace the positive target
$Y_0=1$ by $G(\omega)=1+K(\omega-\omega_c)^2$, $K>0$.
This only replaces the real offdiagonal target by
\[
 b(\omega)\approx
 \frac{a_0\sqrt{G(\omega)}}{2\sin(\tau\omega/2)},\qquad
 \tau=\frac{2\pi}{\inf J+\sup J}.
\]
The signed-resolvent approximation and positive-semidefinite residue
completion are unchanged. The limiting normalized spectrum is
\begin{equation}\label{eq:ideal}
 Y_0=4(b/a)^2\sin^2(\tau\omega a/(2a_0)),\qquad \Gamma_0=1.
\end{equation}
Choose the total uniform error, including the finite-stage and damping
perturbations, smaller than one quarter of
$\min_{J\setminus U}[G-G(\omega_c)]$. The finite weak/thick stage and positive
active damping can then be fixed so that its actual
$Y_\infty(\omega,0)$ is positive, all its global minimizers on $J$ lie in $U$,
and its contrast exceeds $1-\epsilon/8$ throughout $J$.
Retain strict upper material-cap margins and strict positive real and
imaginary parts in the active material spaces.

Choose any one actual minimizer $\nu\in U$, put $\lambda=\nu/\omega_c$, and
scale the planar thicknesses by $\lambda$ and its complete material laws by
$\omega\mapsto\lambda\omega$.
Planar power densities satisfy
\[
 Y_{\infty,{\rm new}}(\omega,\theta)
   =Y_{\infty,{\rm old}}(\lambda\omega,\theta),
 \qquad
 \gamma_{\infty,{\rm new}}(\omega,\theta)
   =\gamma_{\infty,{\rm old}}(\lambda\omega,\theta).
\]
There is no area factor for these per-area quantities. Henceforth freeze this
scaled stage and write
\begin{equation}\label{eq:baseline-min}
 Y_c:=Y_\infty(\omega_c,0)>0,\qquad
 Y_\infty(\omega,0)\ge Y_c\quad(\omega\in B).
\end{equation}
No choice of a minimizer as a function of a later tuning parameter occurs.
Several baseline minima cause no difficulty.

Scaling preserves the Lorentz class: a term
$R/(\Omega^2-z^2-i\gamma z)$ evaluated at $\lambda z$ has residue
$R/\lambda^2$, resonance $\Omega/\lambda$ and damping $\gamma/\lambda$.

\begin{lemma}[A passive direction with nonzero real sensitivity]\label{lem:sensitivity}
The stage can be selected so that a real rank-one positive semidefinite
$2\times2$ matrix $T$ in the plate has
\[
 \left.\frac{\partial}{\partial x}
 Y_\infty(\omega_c,0;\zeta=x)\right|_{x=0}\ne0.
\]
Here the plate's active matrix $C$ is changed to $C+\zeta T$;
the physical change is $N^{-1}D_\theta(\zeta T)D_\theta$,
where $D_\theta=\operatorname{diag}(1,\theta)$.
\end{lemma}
\begin{proof}
In \eqref{eq:ideal}, with $a,c$ fixed and $b$ the shared offdiagonal entry,
$\partial_bY_0=2Y_0/b>0$ uniformly on the compact enlarged band.
The exact slab amplitudes and the corresponding normalized complex forward
work are holomorphic in a small complex variation of $b$. The weak/thick
estimates are locally uniform on that disk, as follows from the same bounded
matrix exponentials and small interface reflections used in \cite{strict}.
Apply Cauchy's formula to the complex work, then take the imaginary part of
its real-$b$ derivative. This transfers the displayed extinction derivative
to a sufficiently large finite weak/thick stage; $Y$ itself is not asserted
to be holomorphic. Small positive damping preserves it by finite-stage
continuity. This is a material derivative, not a frequency derivative.

Let $T_\pm=(e_x\pm e_y)(e_x\pm e_y)^T$. Their directional derivatives differ by
$2\partial_bY$, which is nonzero. At least one is therefore nonzero at the
chosen $\nu$. Choose that matrix for $T$ and perform the planar scaling above.
The sensitivity and its nonzero sign are preserved.
\end{proof}

\section{Uniform analyticity in a material parameter}
We state explicitly the strengthening of the old aperture theorem that is
needed. Only the material parameter is differentiated.

Write $\mathcal H_R$ for the absorber scalar and plate two-component auxiliary
current space, and let
\[
 \alpha(\omega,\zeta)
 =\operatorname{diag}\!\left(\alpha_A(\omega),
           [C(\omega)+\zeta T]/N\right),\qquad
 L_{R,\theta}(\omega,\zeta)=
 \alpha(\omega,\zeta)^{-1}-J_\theta^*K_{D_R}(k)J_\theta .
\]
The scaled stage's thicknesses and laws are understood in this notation.
Choose a disk $\abs\zeta<r$ so small that, throughout $B$ and its closure,
$\ReH\alpha$ and $\ImH\alpha$ remain strictly positive on the active spaces,
and the physical upper caps retain margins for a small fixed interval
$0\le\theta\le\theta_0$.
This is possible by compactness and the frozen positive active losses.
The disk includes slight negative imaginary perturbations as an analytic
device; every actual added oscillator will be passive.

As in \cite{strict}, put
\[
 W=-\ImH\alpha^{-1}\succeq c_*I,\qquad
 -\Im\ip r{L_{R,\theta}r}\ge c_*\norm r^2 ,
\]
uniformly in the compact parameters and independently of $R$.
Define the \emph{complex} normalized work
\[
 \mathcal F_R(\omega,\theta,\zeta)
 =kR^{-2}\ip{g_+}{L_{R,\theta}(\omega,\zeta)^{-1}g_+},
 \qquad
 \Im\mathcal F_R=R^{-2}Y_R.
\]
The fixed incident vector $g_+$ does not depend on $\zeta$, so this expression
is holomorphic in $\zeta$.

\begin{lemma}[Material-parameter aperture transfer]\label{lem:cauchy}
On compact subdisks and on the fixed real frequency and coupling intervals,
\[
 \partial_\zeta^j\mathcal F_R
   \longrightarrow\partial_\zeta^j\mathcal F_\infty
 \quad\text{uniformly},\qquad j=0,1,2,\ldots .
\]
The normalized absorption and contrast converge uniformly in value wherever
the planar extinction is positive.
\end{lemma}
\begin{proof}
The current-localization proof of \cite{strict} is uniform when the continuous
active matrices range over this compact disk: their inverse coercivity
constant, weights and height norms have common bounds, and the same Weyl
estimates apply. Its $o(R)$ current error and the $O(R)$ incident norm imply
convergence of the full complex work, not just its imaginary part.
Apply Cauchy's integral formula on a slightly larger subdisk to obtain every
fixed material derivative. The absorption and contrast statement uses the
original quadratic work estimates and positive denominator.
\end{proof}

In particular, for $s$ in any fixed positive compact interval and
$\theta=s/R$, the limiting complex work is
$\mathcal F_\infty(\omega,0,\zeta)$, with the same material-derivative
convergence. No derivative in $\omega$ is asserted.

\section{A weak passive oscillator creates an extinction dip}
At the frozen center let
\[
 F(\zeta)=\mathcal F_\infty(\omega_c,0,\zeta),\qquad
 F'(0)=A+iB_1 .
\]
Here $F$ includes the factor $k_c$; Lemma~\ref{lem:sensitivity} gives
$B_1=\partial_{\Re\zeta}Y_\infty\ne0$.
Put
\[
 \ell(u)=\frac1{-u-i}=\frac{-u+i}{1+u^2},\qquad
 \Phi_\rho(u)=\Im F(\rho\ell(u)).
\]
For small $\rho>0$,
\begin{equation}\label{eq:fano}
 \Phi_\rho(u)=Y_c+
       \rho\,\frac{A-B_1u}{1+u^2}+O(\rho^2).
\end{equation}
The error is uniform on the compactified real line and in $C^2$ on bounded
$u$ intervals. The positive imaginary part of $\ell$ will be supplied by an
ordinary passive Lorentz oscillator.

\begin{lemma}[Strict frozen-profile minimum]\label{lem:fano}
For sufficiently small $\rho$, $\Phi_\rho$ has a unique global minimum
at a finite $u_\rho$, with positive second derivative there, and
\[
 m_\rho:=\Phi_\rho(u_\rho)=Y_c-d_\rho,\qquad
 0<d_\rho<Y_c/4.
\]
The minimum has a strict gap away from any sufficiently small neighborhood
of $u_\rho$.
\end{lemma}
\begin{proof}
Let $D=(A^2+B_1^2)^{1/2}$. The leading rational function in
\eqref{eq:fano} has its unique minimum at
\[
 u_*=\frac{A+D}{B_1},\qquad
 h(u_*)=\frac{A-D}{2}<0,\qquad
 h''(u_*)=\frac{2D}{(1+u_*^2)^2}>0.
\]
Its value tends to zero at both ends of the real line. It therefore has a
positive value gap outside a small neighborhood of $u_*$. Uniform $C^0$
control preserves that gap, and $C^2$ control preserves strict convexity
inside the neighborhood. This proves uniqueness, nondegeneracy and
$d_\rho=\rho(D-A)/2+O(\rho^2)>0$. Decrease $\rho$ to obtain the displayed
upper bound and to stay in the analytic disk.
\end{proof}

Choose $\rho$ still smaller if necessary so that all material values
$\abs\zeta\le2\rho$ preserve the cap margins, positive extinction, and planar
contrast greater than $1-\epsilon/4$ on the band and near $\theta=0$.
This use of continuity is legitimate even if $d_\rho$ becomes very small:
the aperture has not yet been chosen. The added resonance reduces extinction
by interference; passivity does not require extinction to increase under every
positive-semidefinite material perturbation.

\section{Choose the finite aperture before narrowing the resonance}
Set
\[
 s_-=\sqrt{\frac{E_*}{2Y_c}},\qquad
 s_+=\sqrt{\frac{2E_*}{Y_c}},\qquad I_s=[s_-,s_+].
\]
Choose a sufficiently large \emph{finite} $R$, so that $\theta=s/R$ is in the
small coupling interval for all $s\in I_s$.
By Lemma~\ref{lem:cauchy}, the frozen finite profiles
\[
 \Phi_{R,s}(u)=R^{-2}Y_R(\omega_c,s/R;\rho\ell(u))
\]
are uniformly close to $\Phi_\rho$ in $C^2$ near its minimum and in $C^0$
on the full profile. They therefore have unique minima in one common compact
$u$ interval, uniformly positive local curvature, and uniform gaps away from
that interval. Moreover the unperturbed whole-band spectra satisfy
\[
 R^{-2}Y_R(\omega,s/R;0)\ge Y_c-\delta_R,\qquad
 \delta_R\longrightarrow0.
\]
Choose $R$ so large that all these errors are much smaller than $d_\rho$.
The contrast is uniformly greater than $1-\epsilon/2$ for the material
parameters under consideration.

At this selected finite aperture add a sufficiently small common isotropic
Lorentz susceptibility $\nu f_0(\omega)I$ in both layers, with positive
residue and damping and resonance above the band.
The electric resolvent formula and Neumann perturbation from \cite{strict}
apply uniformly in $\omega,s,\zeta$ on a compact subdisk. They preserve the
strict caps and contrast margins. For fixed $R$, complex work is still
holomorphic in $\zeta$, so Cauchy's formula preserves the required
material-derivative margins as $\nu\downarrow0$.

The physical extinction no longer factors at zero coupling after this
addition. On the positive bracket define, using the physical polarization
current $p_\nu$,
\begin{equation}\label{eq:finite-F}
 \mathcal W_\nu(\omega,s,\zeta)=\frac{k}{s^2}\ip{f_+}{p_\nu},
 \qquad
 \mathcal Y(\omega,s,\zeta)=\Im\mathcal W_\nu
 =\frac{X^{(\nu)}_R(\omega,s/R;\zeta)}{s^2}.
\end{equation}
The complex work $\mathcal W_\nu$ is holomorphic in $\zeta$ and converges
uniformly as $\nu\downarrow0$ on a slightly larger subdisk. Cauchy's formula
therefore supplies convergence of its material derivatives on the working
disk. Before regularization, $\mathcal Y=R^{-2}Y_R$ exactly. Since $s_-$ is
positive, this normalization is uniformly harmless on the full bracket.
Choose all errors small enough that, uniformly in $s$,
\begin{align}
 \mathcal Y(\omega,s,0)&\ge Y_c-d_\rho/8
       &&(\omega\in B),\label{eq:baseline-gap}\\
 \phi_s(u):=\mathcal Y(\omega_c,s,\rho\ell(u))
 &\text{ has a unique nondegenerate minimum }u_s,\nonumber\\
 \phi_s(u_s)&\le Y_c-7d_\rho/8,\label{eq:dip-gap}\\
 \tfrac34Y_c&<\phi_s(u)<\tfrac54Y_c
       &&\text{near the minima}.\label{eq:value-bracket}
\end{align}
The intervals of strict curvature and the gaps can be chosen uniformly.
The minimum locations $u_s$ depend continuously on $s$.
All original oscillators and the isotropic regularizer already have positive
damping. Nothing has yet required a narrow frequency resonance.

For fixed $R$, the baseline and its complex work are smooth in real frequency
on a neighborhood of $B$, jointly with the finite parameter sets.
Indeed, the static singular part of $K_{D_R}(k)$ is independent of $k$;
its remaining kernel is weakly singular and operator-norm holomorphic on
bounded complex frequency sets. The invertible volume operator then has a
holomorphic inverse locally about each real frequency.
Only this fixed-body regularity is used below.

\section{Put the actual minimum exactly at the center}
For a positive width $\gamma$ and a real location parameter $v$ in a fixed
compact interval about the $u_s$, add to the active plate the Lorentz term
$\zeta_{\gamma,v}(\omega)T/N$, where
\begin{equation}\label{eq:notch}
 \Omega_v^2=\omega_c^2-\gamma\omega_cv,\qquad
 \zeta_{\gamma,v}(\omega)=
 \frac{\rho\gamma\omega_c}
      {\Omega_v^2-\omega^2-i\gamma\omega}.
\end{equation}
For sufficiently small $\gamma$, $\Omega_v^2>0$ uniformly in $v$.
The final physical residue is
$(\rho\gamma\omega_c/N)D_\theta TD_\theta\succeq0$.
Thus this is one ordinary passive reciprocal oscillator with positive damping.
Its peak perturbation is small, although its frequency derivatives may be large.

At the center its material value is \emph{exactly}
\[
 \zeta_{\gamma,v}(\omega_c)=\rho\ell(v).
\]
With $\omega=\omega_c+\gamma t/2$, direct algebra gives
\begin{equation}\label{eq:scaled-lorentz}
 \zeta_{\gamma,v}(\omega_c+\gamma t/2)=
 \frac{\rho}
 {-v-t-\gamma t^2/(4\omega_c)-i(1+\gamma t/(2\omega_c))}
 \longrightarrow \rho\ell(t+v)
\end{equation}
in $C^2(t,v)$ on every fixed compact set. Also
$\sup_{\omega\in B,v}\abs{\zeta_{\gamma,v}(\omega)}
\le(1+o(1))\rho$, so the perturbation remains in the material disk.
For sufficiently large $\abs t$ and $\omega\in B$,
\begin{equation}\label{eq:tail}
 \abs{\zeta_{\gamma,v}(\omega)}\le C_B\rho/\abs t.
\end{equation}
For example the real part of the denominator in \eqref{eq:scaled-lorentz}
is $-v-t(\omega+\omega_c)/(2\omega_c)$, proving this bound for bounded $v$.

At the now-fixed aperture define
\[
 H_\gamma(s,t,v)=
 \mathcal Y(\omega_c+\gamma t/2,s,
               \zeta_{\gamma,v}(\omega_c+\gamma t/2)).
\]
Fixed-body smoothness and \eqref{eq:scaled-lorentz} imply
\begin{equation}\label{eq:inner-limit}
 H_\gamma(s,t,v)\longrightarrow\phi_s(t+v)
 \quad\text{in }C^2(t,v),
\end{equation}
uniformly for $s\in I_s$ and bounded $t,v$.
At $\gamma=0$, the equation $\partial_tH_0(s,0,v)=0$ has the unique relevant
root $v=u_s$, with $\partial_v\partial_tH_0=\phi_s''(u_s)>0$.
The uniform implicit-function argument therefore supplies a continuous
$v_\gamma(s)$ for all sufficiently small fixed $\gamma>0$, with
\begin{equation}\label{eq:stationary}
 \partial_\omega X_R^{(\nu)}(
   \omega_c,s/R;\zeta_{\gamma,v_\gamma(s)}(\omega))=0
\end{equation}
where the derivative is of the complete composed frequency response.
The local curvature is positive. One may also obtain the root directly by
strict monotonicity and opposite endpoint signs of
$\partial_tH_\gamma(s,0,v)$.

We verify \emph{global}, not just local, minimality.
First choose a finite $T$ sufficiently large that the tail bound
\eqref{eq:tail} and the fixed-body material Lipschitz bound change
$\mathcal Y$ by less than $d_\rho/8$ whenever $\abs t\ge T$.
This choice can be made after $R$; it introduces no aperture-uniform
frequency estimate. By \eqref{eq:baseline-gap}, every such outer value is at
least $Y_c-d_\rho/4$, while the center remains below $Y_c-3d_\rho/4$.
For the bounded inner interval choose $a,\mu>0$ so that
$\phi_s''\ge\mu$ on $\abs{u-u_s}\le2a$, and a uniform positive value gap
on $\abs{u-u_s}\ge a/2$. Take $\gamma$ so small that
$\abs{v_\gamma(s)-u_s}<a/2$ and
$\partial_t^2H_\gamma\ge\mu/2$ on $\abs t\le a$.
Stationarity at zero then gives strict minimality on that interval.
For $a\le\abs t\le T$, the frozen profile gap and uniform $C^0$ convergence
in \eqref{eq:inner-limit} give $H_\gamma(s,t,v_\gamma(s))>
H_\gamma(s,0,v_\gamma(s))$.
Decrease $\gamma$ also so that the inner interval lies in $B$.
These two regions cover the band. Hence, for every $s\in I_s$,
\begin{equation}\label{eq:exact-min}
 X(\omega,s)>X(\omega_c,s)\quad
      (\omega\in B\setminus\{\omega_c\}).
\end{equation}

Finally the center values depend continuously on $s$. By
\eqref{eq:value-bracket} and a sufficiently small choice of $\gamma$,
\[
 X(\omega_c,s_-)<\tfrac58E_*<E_*,
 \qquad
 X(\omega_c,s_+)>\tfrac32E_*>E_*.
\]
One intermediate-value argument supplies $s_*\in I_s$ with
$X(\omega_c,s_*)=E_*$ exactly.
Take the single resulting $\theta=s_*/R$, resonance location
$v_\gamma(s_*)$, geometry and complete material laws.
Equation~\eqref{eq:exact-min}, the preserved caps, and the uniform contrast
margin prove the proposed theorem.
The strict optical-theorem argument of \cite{base} gives $\Gamma<1$ for every
selected compact body, hence nonattainment.

\section{Parameter order and scope of verification}
The dependency order is:
\[
 \begin{gathered}
 \text{guarded planar bowl and finite weak/thick stage}\\
 \longrightarrow\text{one planar minimum alignment and freezing}\\
 \longrightarrow\text{material direction and small notch peak }\rho\\
 \longrightarrow\text{large finite aperture }R\\
 \longrightarrow\text{small isotropic regularization }\nu\\
 \longrightarrow\text{finite tail window }T
 \longrightarrow\text{small positive linewidth }\gamma\\
 \longrightarrow\text{continuous resonance centering}
 \longrightarrow\text{exact amplitude tuning}.
 \end{gathered}
\]
The very narrow linewidth is an additional unrestricted resource.
The proof does not provide a practical size, loss or linewidth bound.
It does not establish the optimum with fixed resources or an isotropic
realization. Material oscillator axes, resonance parameters and geometry are
fixed after the final selections; the frequency does not choose another device.

Three GPT-6.1-sol agents separately examined the construction, the
material-parameter aperture argument and the minimum/global-tail logic.
Their reviews reported no substantive blocker after the stated corrections.
This is AI-assisted analytical checking, not independent human review.
The accompanying \texttt{boundary\_checks.py} and
\texttt{numerical\_results.json} give a reproducible NumPy diagnostic of the
selection argument. The model is an \emph{analytic surrogate}, consisting of
a smooth bowl, fixed small rapid ripples, and a holomorphic quadratic work
response to the added Lorentz susceptibility. It is not asserted to be a
realizable Maxwell response, and no absorption-asymmetry ratio is computed.

The diagnostic checks the Fano minimum and curvature formulas, the exact
center material identity, and the scaled Lorentz derivatives through order
two. Their maximum sampled error decreases from $3.70\times10^{-3}$ at
$\gamma=0.03$ to $3.52\times10^{-6}$ at $\gamma=3\times10^{-5}$.
The broad centered resonance fails the sampled global-minimum check in the
presence of the fixed ripples. At the narrow selected width, nested centering
and amplitude tuning reach the surrogate target $E_*=1.7$ with magnitude
residual below $1.1\times10^{-14}$ and frequency-derivative residual below
$2.2\times10^{-12}$. The center is the sampled minimum on a 259,885-point
hybrid grid, and all 13 sampled amplitude-bracket values pass the analogous
check. These are numerical diagnostics, not a continuum-band certificate or
a finite three-dimensional device calculation.

\section{A separate obstruction to exactly constant extinction}
The boundary theorem requires a minimum, not a constant spectrum.
In fact exactly constant positive extinction on an open band is excluded
for the finite, positively damped Lorentz subclass constructed here.

\begin{proposition}[No positive flat interval in this Lorentz subclass]
Let a fixed compact body have finitely many homogeneous regions and, in each,
a finite oscillator sum
\[
 \chi_j(z)=\sum_l
 \frac{R_{jl}}{\Omega_{jl}^2-z^2-i\gamma_{jl}z},\qquad
 R_{jl}\succeq0,\quad
 \Omega_{jl},\gamma_{jl}>0,\quad
 \sum_lR_{jl}\succ0.
\]
Its plane-wave extinction cannot equal a positive constant on any open
frequency interval.
\end{proposition}
\begin{proof}
For every real $\omega>0$, $\ImH\chi_j(\omega)\succ0$, so $\chi_j$ is
invertible. With $L(\omega)=\chi(\omega)^{-1}-K_D(\omega/c_{\rm v})$,
the radiation identity gives
$-\Im\ip p{L(\omega)p}\ge c_\omega\norm p^2$ for some $c_\omega>0$.
The finitely many regions give a common constant at each fixed frequency.
Complex Lax--Milgram implies bounded invertibility.

On a bounded domain $K_D(z)$ is operator-norm holomorphic: its static
singular part is fixed, and its remaining kernel and all local complex
frequency derivatives have integrable weak singularities.
Local inversion therefore makes the forward complex work
\[
 W(z)=\int_D e_y\ee^{-izx_3/c_{\rm v}}\cdot
 [L(z)^{-1}(e_y\ee^{izx_3/c_{\rm v}})]\,\dd x
\]
holomorphic in a neighborhood of each positive real frequency.
Thus $X(\omega)=(\omega/c_{\rm v})\Im W(\omega)$ is real analytic on
the connected interval $(0,\infty)$.

At zero, $\chi(0)\succ0$ and the static volume operator is self-adjoint
nonpositive: its whole-space Fourier symbol is
$-\xi\xi^T/\abs\xi^2$ after extension by zero.
Consequently $L(0)=\chi(0)^{-1}-K_D(0)$ is strictly positive and invertible.
The current and $W$ remain bounded near zero, so $X(\omega)\to0$.
Indeed $W(0)$ is real and $X=O(\omega^2)$.
If $X$ were a positive constant on an open interval, the real-analytic
identity theorem would make it that constant throughout $(0,\infty)$,
contradicting this limit.
\end{proof}

This proposition concerns the stated finite Lorentz subclass. It is not a
claim about every possible passive law with nonanalytic boundary values.
It sharpens the stronger-flat-spectrum caveat in the earlier working note;
it places no obstruction on the isolated center minimum constructed here.

\begin{thebibliography}{9}
\bibitem{base}
D. Wang and J. Zheng,
\emph{Unit Supremum of Absorption Asymmetry at Prescribed Extinction for
Compact Reciprocal Dielectrics}, manuscript dated 17 October 2026,
source revision \texttt{22fed75}.
\url{https://dan-k391.github.io/oq7-papers/prescribed-extinction/}.
\bibitem{strict}
\emph{Uniform absorption asymmetry on a fixed frequency band:
A constructive draft for a strict extinction floor},
research working note, 17 October 2026, Sections 2--6.
Local source in the sibling directory
\texttt{broadband-fixed-band-2026-10-17/main.tex}.
Source SHA-256:
\texttt{1947dd3a23343d78146415c2760a1ac53699f1caa0a0fc44b79af55c560d8f10}.
\end{thebibliography}
\end{document}
